\documentclass[aspectratio=169,xcolor={dvipsnames,table}]{beamer}
\usepackage[no-math,deluxe,haranoaji]{luatexja-preset}
\renewcommand{\kanjifamilydefault}{\gtdefault}
\renewcommand{\emph}[1]{{\upshape\bfseries #1}}
\usetheme{metropolis}
\metroset{block=fill}
\setbeamertemplate{navigation symbols}{}
\setbeamertemplate{blocks}[rounded][shadow=false]
\usecolortheme[rgb={0.7,0.2,0.2}]{structure}
%%%%%%%%%%%%%%%%%%%%%%%%%%
%% Change alert block colors
%%% 1- Block title (background and text)
\setbeamercolor{block title alerted}{fg=mDarkTeal, bg=mLightBrown!45!yellow!45}

\setbeamercolor{block title example}{fg=magenta!10!black, bg=mLightGreen!70}
%%% 2- Block body (background)
\setbeamercolor{block body alerted}{bg=mLightBrown!25}
\setbeamercolor{block body example}{bg=mLightGreen!15}
%%%%%%%%%%%%%%%%%%%%%%%%%%%
\usepackage[absolute,overlay]{textpos}
%\usepackage[grid=true,gridcolor=Maroon,subgridcolor=gray,gridunit=pt,texcoord]{eso-pic} %場所決めのためのgrid表示
%%%%%%%%%%%%%%%%%%%%%%%%%%%
%% さまざまなアイコン
%%%%%%%%%%%%%%%%%%%%%%%%%%%
%\usepackage{fontawesome}
\usepackage{fontawesome5}
\usepackage{figchild}
\usepackage{twemojis}
\usepackage{utfsym}
\usepackage{bclogo}
\usepackage{marvosym}
\usepackage{fontmfizz}
\usepackage{pifont}
\usepackage{phaistos}
\usepackage{worldflags}
\usepackage{jigsaw}
\usepackage{tikzlings}
\usepackage{tikzducks}
\usepackage{scsnowman}
\usepackage{epsdice}
\usepackage{halloweenmath}
\usepackage{svrsymbols}
\usepackage{countriesofeurope}
\usepackage{tipa}
\usepackage{manfnt}
%%%%%%%%%%%%%%%%%%%%%%%%%%%
\usepackage{tikz}
\usetikzlibrary{calc,patterns,decorations.pathmorphing,backgrounds}
\usepackage{tcolorbox}
\usepackage{tikzpeople}
\usepackage{circledsteps}
\usepackage{xcolor}
\usepackage{amsmath}
\usepackage{booktabs}
\usepackage{chronology}
\usepackage{signchart}
%%%%%%%%%%%%%%%%%%%%%%%%%%%
%% 場合分け
%%%%%%%%%%%%%%%%%%%%%%%%%%%
\usepackage{cases}
%%%%%%%%%%%%%%%%%%%%%%%%%%
\usepackage{pdfpages}
%%%%%%%%%%%%%%%%%%%%%%%%%%%
%% 音声リンク表示
\newcommand{\myaudio}[1]{\href{#1}{\faVolumeUp}}
%%%%%%%%%%%%%%%%%%%%%%%%%%
%% \myAnch{<名前>}{<色>}{<テキスト>}
%% 指定のテキストを指定の色の四角枠で囲み, 指定の名前をもつTikZの
%% ノードとして出力する. 図には remember picture 属性を付けている
%% ので外部から参照可能である.
\newcommand*{\myAnch}[3]{%
  \tikz[remember picture,baseline=(#1.base)]
    \node[draw,rectangle,line width=1pt,#2] (#1) {\normalcolor #3};
}
%%%%%%%%%%%%%%%%%%%%%%%%%%
%% \myEmph コマンドの定義
%%%%%%%%%%%%%%%%%%%%%%%%%%
%\newcommand{\myEmph}[3]{%
%    \textbf<#1>{\color<#1>{#2}{#3}}%
%}
\usepackage{xparse} % xparseパッケージの読み込み
\NewDocumentCommand{\myEmph}{O{} m m}{%
    \def\argOne{#1}%
    \ifx\argOne\empty
        \textbf{\color{#2}{#3}}% オプション引数が省略された場合
    \else
        \textbf<#1>{\color<#1>{#2}{#3}}% オプション引数が指定された場合
    \fi
}
%%%%%%%%%%%%%%%%%%%%%%%%%%%
%%%%%%%%%%%%%%%%%%%%%%%%%%%
%% 文末の上昇イントネーション記号 \myRisingPitch
%% 通常のイントネーション \myDownwardPitch
%% https://note.com/dan_oyama/n/n8be58e8797b2
%%%%%%%%%%%%%%%%%%%%%%%%%%%
\newcommand{\myRisingPitch}{
\begin{tikzpicture}[scale=0.3,baseline=0.3]
\draw[->,>=stealth] (0,0) to[bend right=45] (1,1);
\end{tikzpicture}
}
\newcommand{\myDownwardPitch}{
\begin{tikzpicture}[scale=0.3,baseline=0.3]
\draw[->,>=stealth] (0,1) to[bend left=45] (1,0);
\end{tikzpicture}
}
%%%%%%%%%%%%%%%%%%%%%%%%%%%%
%\AtBeginSection[%
%]{%
%  \begin{frame}[plain]\frametitle{授業の流れ}
%     \tableofcontents[currentsection]
%   \end{frame}%
%}

\usepackage{pxrubrica}
\usetikzlibrary{tikzmark}
%%%%%%%%%%%%%%%%%%%%%%%%%%%
\title{English is fun.}
\subtitle{She studied hard to pass the exam.}
\author{}
\institute[]{}
\date[]

%%%%%%%%%%%%%%%%%%%%%%%%%%%%
%% TEXT
%%%%%%%%%%%%%%%%%%%%%%%%%%%%
\begin{document}

\begin{frame}[plain]
  \titlepage
\end{frame}

\section*{授業の流れ}
\begin{frame}[plain]
  \frametitle{授業の流れ}
  \tableofcontents
\end{frame}

\section{副詞的用法}
\subsection{目的}
%%%%%%%%%%%%%%%%%%%%%%%%%%%%%%%%%%%%%%%%%%%%%
\begin{frame}[plain]{～するために}
 \large

\begin{enumerate}
 \item<1-> You must work hard.\hfill{\scriptsize work \textipa{/w\'\textrhookschwa :k/} 働く、勉強する}
 \item<2-> You must \myAnch{a}{white}{work}hard \,\,\,\myAnch{b}{Maroon}{\textbf{to pass} the exam}.\hfill{\scriptsize pass \textipa{/p\'\ae s/} （試験に）受かる}

\end{enumerate}

\begin{tikzpicture}[remember picture, overlay]
% Calculate intermediate points
  \coordinate (a1) at ($(a) + (0,-25pt)$); % 25pt below a
  \coordinate (b1) at ($(b) + (0,-25pt)$); % 25pt below b
  % Draw the arrow with right angles
  \visible<3->{\draw[<-,Maroon,line width=2pt, opacity=.5] (a) -- (a1) -- (b1) -- (b);}
\end{tikzpicture}

%
\hfill{\tiny 0113}\,{\scriptsize \myaudio{./audio/034_infinitive_adv_01.mp3}}
%
\begin{block}<4->{Topic for Today}
\begin{itemize}\setbeamertemplate{items}[square]\small
 \item to不定詞が「〜するために」の意味で動詞を修飾することがある
 \item \kenten{目的}を表します
 \item to不定詞の\kenten{副詞的用法}といいます\\
       \hfill{\scriptsize 動詞を修飾するのは副詞の役目ですよね}
 \end{itemize}
     \end{block}

\end{frame}
%%%%%%%%%%%%%%%%%%%%%%%%%%%%%%%%%%%%%%%%%%%%%
\section{目的を表すto不定詞の位置}
%%%%%%%%%%%%%%%%%%%%%%%%%%%%%%%%%%%%%%%%%%%%%
\begin{frame}[plain,t]{目的を表すto不定詞の位置}

 \begin{enumerate}
  \item<2-> \begin{enumerate}
	 \item<2-> The door opened \textbf{suddenly}%
	      \hfill\visible<3->{{\scriptsize S $+$ V\,\,\,\Circled{\,副詞\,}}}
	 \item<4-> \textbf{Suddenly}, the door opened. 
	      \hfill\visible<5->{{\scriptsize \Circled{\,副詞\,},\,\,\,S $+$ V }}
	\end{enumerate}
       \hfill{\scriptsize suddenly \textipa{/s\'\textturnv dnli/ 突然}}%
  \item<6-> \begin{enumerate}
	 \item<6-> You must work hard \fbox{\textbf{to pass} the exam}.%
	      \hfill\visible<7->{{\scriptsize S $+$ V\,\,\,\Circled{\,to不定詞\,}}}
	 \item<8-> \fbox{\textbf{To pass} the exam}, you must work hard.%
	      \hfill\visible<9->{{\scriptsize \Circled{\,To不定詞\,},\,\,\,S $+$ V }}
	\end{enumerate}
 \end{enumerate}

\vspace{20pt}

\begin{block}<10->{Topic for Today}
\begin{itemize}\setbeamertemplate{items}[square]\small
 \item \kenten{目的}をあらわす副詞的用法のto不定詞は、文の最後にも先頭にもきます
       \begin{itemize}\setbeamertemplate{items}[circle]\small
\setlength{\itemsep}{5pt} % 項目間
	\item $\text{
S} + \text{V}\,\,\,\,$\Circled[fill color= yellow!40]{\,to不定詞\,}
	\item \Circled[fill color= yellow!30]{\,To不定詞\,},\,\, $\text{S} + \text{V}$.
       \end{itemize}
 \end{itemize}
\hfill{\scriptsize1語の副詞の場合と同じです}
     \end{block}

\vspace{-10pt}

\hfill{\tiny 0218}\,{\scriptsize \myaudio{./audio/034_infinitive_adv_02.mp3}}

\end{frame}
%%%%%%%%%%%%%%%%%%%%%%%%%%%%%%%%%%%
\begin{frame}[plain]{Exercises}

{\small 各英文の意味を考えましょう}

 \begin{enumerate}
  \item I got up early \textbf{to catch} the first train.\hfill{\scriptsize catch \textipa{/k\'\ae\textteshlig/} 間に合う}
  \item She went to the station \textbf{to meet} her parents.\hfill{\scriptsize meet \textipa{/m\'\i:t/} 出迎える}
  \item \textbf{To save} money, we eat at home.\hfill{\scriptsize save \textipa{/s\'eIv/ 節約する}}
  \item \textbf{To help} his parents, he cleans the house every day.\hfill{\scriptsize clean \textipa{/kl\'\i:n/ 掃除する}}
 \end{enumerate}

\vspace{100pt}
 
\hfill{\tiny 0218}\,{\scriptsize \myaudio{./audio/034_infinitive_adv_02.mp3}}

\normalsize
\begin{textblock*}{0.4\linewidth}(310pt,137pt)
\visible<2->{\begin{tikzpicture}
\pig[
signpost=\scalebox{0.5}{
\parbox{2.2cm}{\color{black}
\centering 文の先頭にも\\最後にも\\きますね！}},
speech={\tiny
 \begin{tabular}{c}
 目的を表す\\to不定詞
 \end{tabular}
},
signcolour= brown!70!gray,
signback=white!80!brown
]
\end{tikzpicture}}
\end{textblock*}

\end{frame}
%%%%%%%%%%%%%%%%%%%%%%%%%%%%%%%%%%%
\begin{frame}[plain]{Exercises}

{\small 日本語の意味になるよう(~~~~~~)内の語句を並べ替えましょう。先頭の語は大文字で始めてください}%
\hfill{\tiny 0212}\,{\scriptsize \myaudio{./audio/034_infinitive_adv_03.mp3}}
 \begin{enumerate}
  \item {\small 彼らは野球をするために公園へ行った。}\\
	They ( play / to / to / the park / went / baseball ) .\\
	\visible<2->{They went to the park to play baseball.}
  \item {\small 彼はフランスへ行くためにフランス語を勉強した。}\\
	He ( to / to / French / studied / go / France ) .\\
	\visible<3->{He studied French to go to France.}
  \item {\small 彼女は新しい自転車を買うためにお金を節約した。}\\
	To ( saved / bike / money / new / she / buy / a ) .\\
	\visible<4->{To buy a new bike, she saved money.}
  \item {\small そのバスに間に合うようにわたしたちは走った。}
	\hfill{}( bus / catch / to / the ), we ran.\\
	\visible<5->{To catch the bus, we ran.}
 \end{enumerate}
\end{frame}
%%%%%%%%%%%%%%%%%%%%%%%%%%%%%%%%%%%%%
\section{To不定詞で文がはじまったとき}
%%%%%%%%%%%%%%%%%%%%%%%%%%%%%%%%%%%%
\begin{frame}[plain]{to不定詞で文がはじまったとき}

 {\small つぎの2文の意味を考えましょう\textdbend}

\begin{enumerate}
 \item \alt<2->{\fbox{\textbf{To study} English}}{To study English}\,\,\,is very important.
\hfill{}%
\visible<4->{\Circled[fill color= yellow!30]{\,To不定詞($=$ 名詞的用法)\,}\, ($=$S)\,\, $+$ V}

 \item \alt<2->{\fbox{\textbf{To study} English}}{To study English}\,,\,\,\,he went to Australia.\hfill%
\visible<5->{\Circled[fill color= yellow!30]{\,To不定詞($=$ 副詞的用法)\,},\,\,\,\, $\text{S} + \text {V}$}\\
\hfill{\scriptsize Australia \textipa{/O:str\'eIlj@/} オーストラリア}
\end{enumerate}
%
\hfill{\tiny 0120}\,{\scriptsize \myaudio{./audio/034_infinitive_adv_04.mp3}}


\bigskip 

\hfill\begin{minipage}{.4\textwidth}
\begin{itemize}\small
 \item<3->[hint 1] 主語と述語動詞を確認しよう
 \item<3->[hint 2] to不定詞は何的用法か考えよう
\end{itemize}
\end{minipage}
\end{frame}
%%%%%%%%%%%%%%%%%%%%%%%%%%%%%%%%%
\section{to不定詞のまとめ}
%%%%%%%%%%%%%%%%%%%%%%%%%%%%%%%%%%%%%%%%%%%%%
\begin{frame}[plain]{to不定詞とは}
\small
\visible<2->{{\small 動詞とは本来、主語と組み合わせて使うもの}
\mbox{}\hfill{}\myEmph[2-]{NavyBlue}{I} \myEmph[2-]{Maroon}{play}\tikzmark{play} the piano and \myEmph[2-]{NavyBlue}{she} \myEmph[2-]{Maroon}{sings}\tikzmark{sings}.}

\visible<3->{
\hfill{\scriptsize 主語と組み合わせる動詞を特に\Circled{述語動詞}\tikzmark{predicate}といいます}
}


\visible<5->{だが、特別に

\begin{tcolorbox}[colback=yellow!10!white]
\myAnch{v}{white}{動詞}%
\hspace{25pt}%
\begin{tabular}{ll}
 \visible<5->{\myAnch{n}{white}{名詞のはたらき}}&\visible<6->{\fcolorbox{black}{white}{\textbf{To get} up early} is important.}\\
 \visible<5->{\myAnch{adj}{white}{形容詞のはたらき}}&\visible<7->{I want something \fcolorbox{black}{white}{\textbf{to drink}}.}\\
 \visible<5->{\myAnch{adv}{white}{副詞のはたらき}}&\visible<8->{She studied hard \fcolorbox{black}{white}{\textbf{to pass} the exam}}.
\end{tabular}
\end{tcolorbox}

をすることがある
}

\begin{tikzpicture}[remember picture, overlay]
 \visible<5->{\draw[line width=2pt,opacity=.75, gray, ->] (v.east) to[out=0, in=180] (n.west);} 
 \visible<5->{\draw[line width=2pt,opacity=.75, gray, ->] (v.east) to[out=0, in=180] (adj.west);} 
 \visible<5->{\draw[line width=2pt,opacity=.75, gray, ->] (v.east) to[out=0, in=180] (adv.west);} 
\end{tikzpicture}

\visible<9->{
\Circled[fill color = yellow!10!white]{\,$\text{to} + \text{動詞の原形}$\,}
のことをto\kenten{不定詞}といいます
}

\visible<10->{{\scriptsize \textdbend\,\,toは\kenten{述語動詞ではない}という目印！}}%
\hfill{\tiny 0140}\,{\scriptsize \myaudio{./audio/031_infinitive_intro_01.mp3}}

\begin{tikzpicture}[remember picture,overlay]
 \visible<4->{\draw[Stealth-,opacity=.4,line width=2pt] ([yshift=-2pt]pic cs:play) to[out=-30,in=160] ([xshift=-16pt, yshift=10pt] pic cs:predicate);}
 \visible<4->{\draw[Stealth-,opacity=.4,line width=2pt] ([xshift=-25pt,yshift=-1pt]pic cs:sings) to ([xshift=-16pt, yshift=10pt] pic cs:predicate);}
\end{tikzpicture}

\begin{textblock*}{0.4\linewidth}(300pt,170pt)
\visible<9->{\begin{tikzpicture}
\duck[signpost=\scalebox{0.3}{
\parbox{2.5cm}{\color{black}
{\Large to $+$ 原形}\\{\Large $=$ to不定詞}}},
signcolour=brown!70!gray,
signback=white!80!brown,
graduate=gray!20!black,
tassel=red!70!black,
speech={\scalebox{.75}{\tiny 述語動詞じゃない}}
]
\end{tikzpicture}}
\end{textblock*}

\end{frame}
%%%%%%%%%%%%%%%%%%%%%%%%
\end{document}

